\begin{lemma}[Exit vertices miss no free color; Scheide Lemma 3.8]
For every globally maximum triple $(T,e_0,\varphi)$,
\[
 \overline\varphi(F(T,e_0,\varphi))
 \cap\mathcal F(T,e_0,\varphi)=\varnothing.
\]
Equivalently, no Scheide exit vertex misses a free color.
\end{lemma}
\begin{proof}
Let $v\in F(T,e_0,\varphi)$.  By
\noderef{ScheideExitVertex}, there are a tree vertex $u$, a color
$\gamma$ missing at $u$, a defective color $\delta$, and a full
$(\delta,\gamma)$-path $P$ from $u$ to an outside vertex $z$, with $v$ the
last tree vertex on $P$.  Suppose that a free color $\alpha$ is missing at
$v$.

If $\alpha=\gamma$, elementariness from
\noderef{ScheideGlobalMaximumOrder} forces $u=v$.  Apply
\noderef{ScheideDefectiveChainDepartureReturn} to the full path $P$; it
produces a first departure and a last later tree vertex $v_2$.  The exit
witness says that $v$ is the last tree vertex on the same directed path,
so $v_2=v$.  Now
\noderef{ScheideDefectiveChainLastReturnNoFree}, with the free color
$\gamma=\alpha$ and defective color $\delta$, says that $v_2=v$ misses no
free color, a contradiction.

Assume $\alpha\ne\gamma$.  Both colors are missing on the closed set
$V(T)$, so \noderef{ScheideGlobalMaximumBoundaryFacts} says neither occurs
on a boundary edge identity.  Interchange $\alpha$ and $\gamma$ on every
edge not incident with $T$, using \noderef{OutsideTreeColorSwap}.  The new
coloring $\varphi'$ agrees with $\varphi$ on every edge incident with
$V(T)$.  Thus it remains proper with only $e_0$ uncolored, the displayed
tree remains Tashkinov, its vertex set is still globally maximum, and its
free and defective color sets are unchanged.

Under $\varphi'$, the terminal subpath $vPz$ is the full
$(\delta,\alpha)$-chain from $v$: the swap changes the outside
$\gamma$-segments to $\alpha$-segments and, because $v$ was the last tree
vertex, this subpath has exactly one boundary crossing.  Apply
\noderef{ScheideDefectiveChainBoundaryEdges} to the globally maximum triple
$(T,e_0,\varphi')$, the defective color $\delta$, and the free color
$\alpha$ missing at $v$.  It says that this full chain contains every
$\delta$-colored boundary edge identity.  But
\noderef{ScheideGlobalMaximumBoundaryFacts} gives at least three such
identities, while the terminal subpath contains exactly one boundary edge.
This contradiction proves that no free color is missing at $v$.  Since
$v$ was arbitrary, the asserted disjointness holds for the entire finite
exit set.  The bridge \noderef{ScheideExitVertexTreeExit} records that
these exact exits may subsequently be used wherever the broader tablet
predicate \noderef{TreeExit} is required.
\end{proof}
